% =============================================================================
% Chapter 10: Model Compression — ML Systems Lecture Slides
% =============================================================================
% This deck covers the three dimensions of model compression:
%   1. Structural optimization (pruning, distillation, NAS)
%   2. Precision optimization (quantization FP32→INT8→INT4)
%   3. Architectural efficiency (operator fusion, sparsity)
%
% IMAGES:
%   - optimization-stack.pdf       - quantization-number-line.pdf
%   - pruning-structured-vs-unstructured.pdf
%   - knowledge-distillation-pipeline.pdf
%   - compression-deployment-spectrum.pdf
%   - qat-training-loop.pdf        - nas-overview.pdf
%   - operator-fusion.pdf          - compression-tradeoff-regions.pdf
%   - cover_model_compression.png
% =============================================================================
\documentclass[aspectratio=169, 12pt]{beamer}
\usepackage{../../assets/beamerthememlsys}

\mlsyssetup{
  volume       = {Volume I},
  chapter      = {Chapter 10},
  logo         = {../../assets/img/logo-mlsysbook.png},
  instlogo     = {../../assets/img/logo-harvard.png},
  chaptertitle = {Model Compression},
}

% --- Fonts ---
\usepackage{fontspec}
\setsansfont{Helvetica Neue}[
  BoldFont={Helvetica Neue Bold},
  ItalicFont={Helvetica Neue Italic},
  BoldItalicFont={Helvetica Neue Bold Italic},
]
\setmonofont{JetBrains Mono}[Scale=0.85]

% --- Packages ---
\usepackage{booktabs}
\usepackage{amsmath}

% --- Image paths ---
\graphicspath{
  {images/}
}

% --- Chapter-specific macros ---
% \Dvol and \Rpeak are defined in the theme

% --- Helper: safe image include ---
\newcommand{\safeimg}[2][width=\textwidth,keepaspectratio]{%
  \IfFileExists{images/#2}{\includegraphics[#1]{#2}}{%
    \IfFileExists{#2}{\includegraphics[#1]{#2}}{%
      \fbox{\parbox[c][2.5cm][c]{0.85\linewidth}{%
        \centering\footnotesize\textcolor{midgray}{[Missing image]}}}%
    }%
  }%
}

% --- Section count for navigation ---
\setcounter{mlsystotalsections}{7}

\title{Model Compression}
\author{Vijay Janapa Reddi}
\institute{Harvard University}
\date{}

\begin{document}

% =============================================================================
% TITLE SLIDE
% =============================================================================
\mlsystitle{Model Compression}{From Benchmark Winner to Production Model}{cover_model_compression.png}

% =============================================================================
% LEARNING OBJECTIVES
% =============================================================================
\begin{frame}{Learning Objectives}
\note{[2 min] Read through objectives. Emphasize this chapter bridges training
(Ch8) and hardware (Ch11). The key theme: trained models are too big for
production --- compression makes them deployable. Ask: ``Has anyone tried to
run a large model on a phone?''}

\footnotesize
\begin{enumerate}\setlength\itemsep{0pt}
  \item Explain the \textbf{three-layer optimization stack}: model representation, numerical precision, hardware implementation
  \item Compare \textbf{quantization} strategies (FP32 $\to$ INT8 $\to$ INT4) for memory, energy, and accuracy
  \item Apply \textbf{pruning} (structured vs.\ unstructured) and quantify the accuracy-sparsity trade-off
  \item Implement \textbf{knowledge distillation} to transfer capabilities from teachers to students
  \item Distinguish \textbf{PTQ vs.\ QAT} and identify when each is appropriate
  \item Analyze \textbf{operator fusion} as a zero-cost optimization
  \item Design integrated compression pipelines under deployment constraints
\end{enumerate}

\end{frame}

% =============================================================================
\section{The Deployment Gap}
% =============================================================================

\begin{frame}{Why Compress? The Math Does Not Work}
\note{[3 min] Open with the core tension. A 7B LLM needs 14 GB in FP16 but a
phone has 8 GB. This is not an optimization opportunity --- it is a hard
constraint. Ask: ``What options do you have when 14 GB must fit in 8 GB?''}

\footnotesize
\begin{columns}[T]
  \begin{column}{0.52\textwidth}
    Capability $\neq$ deployability:

    \vspace{0.1cm}
    \renewcommand{\arraystretch}{1.1}
    \begin{tabular}{@{}lr@{}}
      \toprule
      \textbf{Model (FP16)} & \textbf{Memory} \\
      \midrule
      GPT-3 (175B)    & 350 GB \\
      LLaMA-7B        & 14 GB \\
      ResNet-50        & 100 MB \\
      MobileNetV2      & 14 MB \\
      DS-CNN (KWS)     & 500 KB \\
      \bottomrule
    \end{tabular}

    \vspace{0.1cm}
    \begin{mlsyscard}{errorstroke}
    {\scriptsize Smartphone RAM: \textbf{8 GB} (shared with OS).\\
    Microcontroller: \textbf{384 KB} SRAM.\\
    14 GB cannot fit in 8 GB.}
    \end{mlsyscard}
  \end{column}
  \begin{column}{0.45\textwidth}
    \safeimg[width=\textwidth,height=5.0cm,keepaspectratio]{compression-deployment-spectrum.pdf}
  \end{column}
\end{columns}

\end{frame}

\begin{frame}{The Physics of Quantization}
\note{[3 min] Key insight: in silicon, bits represent energy. Moving data costs
orders of magnitude more than computing on it. The energy table is the
``aha'' moment. Ask: ``If memory access costs 40,000x more than
an integer add, where should we focus optimization?''}

\footnotesize
\begin{columns}[T]
  \begin{column}{0.48\textwidth}
    \textbf{Energy-Movement Invariant:}
    $E_{\text{move}} \gg E_{\text{compute}}$

    \vspace{0.1cm}
    \renewcommand{\arraystretch}{1.1}
    \begin{tabular}{@{}llr@{}}
      \toprule
      \textbf{Operation} & \textbf{Width} & \textbf{Energy} \\
      \midrule
      Integer Add & 8-bit  & 1$\times$ \\
      Float Add   & 32-bit & 30$\times$ \\
      DRAM Read   & 64-bit & \textcolor{errorstroke}{\textbf{40,000$\times$}} \\
      \bottomrule
    \end{tabular}

    \vspace{0.1cm}
    \mlsysconcept{Iron Law:}{Reduces \emph{both} $\Dvol$ and $O$.}
  \end{column}
  \begin{column}{0.48\textwidth}
    \begin{mlsyscard}{computestroke}
    {\scriptsize \textbf{The Quantization Speedup}\\[0.05cm]
    7B LLM on 50 GB/s device:\\
    FP16: 14 GB $\to$ 280 ms/token (3.6 tok/s)\\
    INT4: 3.5 GB $\to$ 70 ms/token (14 tok/s)\\[0.05cm]
    \textcolor{computestroke}{\textbf{4$\times$ speedup}} (bandwidth-bound).}
    \end{mlsyscard}

    \vspace{0.1cm}
    FP32 $\to$ INT8: 4$\times$ memory savings.\\
    INT8 energy: \textbf{20$\times$} less per inference.\\
    Battery: 1 hour $\to$ 20 hours.
  \end{column}
\end{columns}

\end{frame}

% =============================================================================
\section{Optimization Stack}
% =============================================================================

\begin{frame}{The Three-Layer Optimization Stack}
\note{[3 min] The organizing framework for the entire chapter. Three layers,
each addressing a different bottleneck. The key insight is they compound:
2x from pruning + 4x from quantization + 2x from fusion = 16x total, not 8x.
Walk through each layer. If short: just show the diagram and summarize.}

% --- Layout: FULL-WIDTH diagram ---
\centering
\safeimg[width=0.88\textwidth,height=4.8cm,keepaspectratio]{optimization-stack.pdf}

\vspace{0.1cm}
\mlsysinsight{Key:}{Techniques \textbf{compound multiplicatively}. Pruning (2$\times$) + Quantization (4$\times$) + Fusion (2$\times$) = \textbf{16$\times$} total compression.}

\end{frame}

\begin{frame}{The Compression-Accuracy Trade-off}
\note{[3 min] Three regions define the optimization landscape. Region 1 is free
--- always do fusion and BF16 first. Region 2 is the sweet spot where most
production optimizations live. Region 3 is the danger zone where models break.
The ``knee'' is where you stop. Common error: students think all quantization
is Region 2 --- INT4 is actually Region 2-3 boundary.}

% --- Layout: FULL-WIDTH diagram ---
\centering
\safeimg[width=0.92\textwidth,height=5.5cm,keepaspectratio]{compression-tradeoff-regions.pdf}

\end{frame}

% --- ACTIVE LEARNING 1: Predict ---
\begin{frame}{Predict: Which Technique First?}
\note{[2 min] Prediction exercise. Give students 60 seconds. The answer is
operator fusion (Region 1: zero accuracy cost). Students often guess
quantization first because it is more dramatic. The point: free optimizations
first, then trade accuracy for efficiency.}

\centering
\vspace{1.0cm}
{\Large\bfseries Think--Write--Share}

\vspace{0.6cm}
{\large You must deploy ResNet-50 to a mobile phone.\\
The model is 100 MB (FP32) and runs at 8 FPS.\\
You need 30 FPS.\\[0.3cm]
\alert{Which optimization technique do you apply first?}}

\vspace{0.5cm}
{\normalsize Write your answer and reasoning. \textcolor{midgray}{(60 seconds)}}

\end{frame}

% =============================================================================
\section{Pruning}
% =============================================================================

\begin{frame}{Pruning: Remove What Doesn't Matter}
\note{[3 min] Neural networks are heavily overparameterized. ResNet-50 can be
pruned by 90\% with less than 2\% accuracy loss. NP-hard, so we use
magnitude heuristics. Ask: ``Why are networks overparameterized?''}

\footnotesize
\begin{columns}[T]
  \begin{column}{0.52\textwidth}
    Networks carry far more parameters than any task requires.

    \vspace{0.05cm}
    \begin{mlsyscard}{crimson}
    {\scriptsize \textbf{Pruning} removes weights that contribute minimal information.}
    \end{mlsyscard}

    \vspace{0.05cm}
    \textbf{Formal objective:}
    $\min_{\hat{W}} \mathcal{L}(\hat{W}) \;\; \text{s.t.} \;\; \|\hat{W}\|_0 \leq k$

    \vspace{0.05cm}
    {\scriptsize
    NP-hard $\to$ magnitude heuristic:
    $M_{i,j} = \begin{cases} 1 & |W_{i,j}| > \tau \\ 0 & \text{otherwise} \end{cases}$
    }
  \end{column}
  \begin{column}{0.45\textwidth}
    {\scriptsize
    \textbf{MobileNet:} 14 MB $\to$ 2 MB, $<$1\% loss\\
    \textbf{ResNet-50:} 90\% prunable (3M of 25.6M params carry task information)
    }

    \vspace{0.1cm}
    \mlsysalert{Pitfall:}{Pruning does \emph{not} auto-speed inference --- sparse matrices need HW support.}
  \end{column}
\end{columns}

\end{frame}

\begin{frame}{Structured vs.\ Unstructured Pruning}
\note{[3 min] The critical distinction. Unstructured: scattered zeros, hardware
cannot skip. Structured: entire columns removed, hardware accelerates.
NVIDIA 2:4 sparsity is the bridge. Common error: all pruning = speedup.}

% --- Layout: FULL-WIDTH diagram ---
\centering
\safeimg[width=0.92\textwidth,height=4.8cm,keepaspectratio]{pruning-structured-vs-unstructured.pdf}

\vspace{0.05cm}
\footnotesize
\begin{columns}[T]
  \begin{column}{0.30\textwidth}
    \centering
    \textcolor{computestroke}{\textbf{Unstructured}}\\
    {\scriptsize Highest compression; needs sparse kernels}
  \end{column}
  \begin{column}{0.30\textwidth}
    \centering
    \textcolor{routingstroke}{\textbf{N:M Structured}}\\
    {\scriptsize 2:4 on A100: 2$\times$ speedup}
  \end{column}
  \begin{column}{0.30\textwidth}
    \centering
    \textcolor{datastroke}{\textbf{Channel/Layer}}\\
    {\scriptsize Dense sub-network; any hardware}
  \end{column}
\end{columns}

\end{frame}

\begin{frame}{Iterative vs.\ One-Shot Pruning}
\note{[2 min] Iterative pruning removes a small fraction per cycle, fine-tunes,
then repeats. One-shot removes everything at once. Iterative loses 0.4\%,
one-shot loses 5\%. If short: just cover the numbers.}

\footnotesize
\begin{columns}[T]
  \begin{column}{0.52\textwidth}
    \textbf{Iterative:} prune $\to$ fine-tune $\to$ repeat

    \vspace{0.05cm}
    {\scriptsize
    \renewcommand{\arraystretch}{1.1}
    \begin{tabular}{@{}lrr@{}}
      \toprule
      \textbf{Cycle} & \textbf{After Prune} & \textbf{After Tune} \\
      \midrule
      1st (2 ch.) & 0.971 & 0.992 \\
      2nd (2 ch.) & 0.956 & 0.993 \\
      3rd (2 ch.) & 0.967 & \textbf{0.991} \\
      \bottomrule
    \end{tabular}
    }

    \vspace{0.05cm}
    27\% channels removed, only \textbf{0.4\%} accuracy loss.

    \vspace{0.1cm}
    \textbf{One-Shot:} prune all at once\\
    {\scriptsize Same 27\%: 0.995 $\to$ 0.914 $\to$ \textbf{0.943} --- \textcolor{errorstroke}{\textbf{5\% degradation}}.}
  \end{column}
  \begin{column}{0.45\textwidth}
    \begin{mlsyscard}{computestroke}
    {\scriptsize \textbf{Lottery Ticket Hypothesis}\\[0.05cm]
    Dense networks contain sparse subnetworks that match full accuracy from original initialization.\\[0.05cm]
    ResNet-18: 10--20\% of size achieves 93.2\% vs.\ 94.1\%.\\[0.05cm]
    \textcolor{computestroke}{\textbf{Implication:}} 80--90\% of training compute may be wasted.}
    \end{mlsyscard}
  \end{column}
\end{columns}
\mlsyscite{Frankle \& Carlin, ICLR 2019}

\end{frame}

% =============================================================================
\section{Distillation}
% =============================================================================

\begin{frame}{Knowledge Distillation: Learning from a Teacher}
\note{[3 min] Teacher soft labels reveal inter-class relationships hard labels
hide. Temperature T controls dark knowledge transfer. DistilBERT: 97\%
accuracy, 40\% fewer params. Ask: ``Why do soft labels help?''}

% --- Layout: FULL-WIDTH diagram ---
\centering
\safeimg[width=0.88\textwidth,height=3.8cm,keepaspectratio]{knowledge-distillation-pipeline.pdf}

\vspace{0.05cm}
\footnotesize
\begin{columns}[T]
  \begin{column}{0.48\textwidth}
    $\mathcal{L} = \alpha \cdot \text{KL}(p^T_{\text{teacher}}, p^T_{\text{student}}) + (1\!-\!\alpha) \cdot \text{CE}$

    \vspace{0.05cm}
    {\scriptsize Temperature $T > 1$ softens distributions, revealing inter-class similarity.}
  \end{column}
  \begin{column}{0.48\textwidth}
    \begin{mlsyscard}{datastroke}
    {\scriptsize \textbf{DistilBERT}: 97\% teacher accuracy, 40\% fewer params, 60\% faster. Produces \emph{dense} models --- any hardware.}
    \end{mlsyscard}
  \end{column}
\end{columns}
\mlsyscite{Hinton et al., 2015}

\end{frame}

\begin{frame}{Distillation vs.\ Pruning}
\note{[2 min] Pruning starts from existing model, makes it sparse. Distillation
starts fresh, produces dense student. In practice, combine both.}

\footnotesize
\renewcommand{\arraystretch}{1.1}
\begin{tabular}{@{}lll@{}}
  \toprule
  \textbf{Criterion} & \textbf{Distillation} & \textbf{Pruning} \\
  \midrule
  \textbf{Mechanism}  & Train student from teacher   & Remove weights from model \\
  \textbf{Output}     & Dense, compact architecture  & Sparse or smaller dense \\
  \textbf{Hardware}   & Runs on any hardware         & May need sparse kernels \\
  \textbf{Compression}& 2--10$\times$                & Up to 10$\times$ \\
  \textbf{Best for}   & Architecture change          & Keep existing architecture \\
  \bottomrule
\end{tabular}

\vspace{0.15cm}
\mlsysconcept{Conservation of Complexity:}{Pruning $\to$ hardware (sparse support). Distillation $\to$ training (teacher budget). Neither destroys complexity.}

\end{frame}

% =============================================================================
\section{Quantization}
% =============================================================================

\mlsysfocus{The Quantization Rule}{%
FP32 $\to$ INT8: \textbf{4$\times$} memory, \textbf{2--4$\times$} throughput\\[0.2cm]
{\normalsize For bandwidth-bound LLMs: linear speedup}%
}

\begin{frame}{How Quantization Works}
\note{[3 min] Walk through the number line diagram. FP32 has continuous values;
INT8 snaps them to 256 discrete levels. The quantization error is the gap
between the true value and the nearest grid point. The formula: x\_q =
round(x / scale) + zero\_point. Emphasize this is a lossy mapping --- you
cannot recover the original value. Ask: ``What happens when two very
different FP32 values map to the same INT8 level?''}

% --- Layout: FULL-WIDTH diagram ---
\centering
\safeimg[width=0.92\textwidth,height=5.0cm,keepaspectratio]{quantization-number-line.pdf}

\vspace{0.05cm}
\small
{\footnotesize Most models exhibit a ``Free Lunch Zone'' where FP32 $\to$ INT8 loses $<$1\% accuracy.\\
Below 4 bits: the \textbf{Quantization Cliff} --- accuracy collapses catastrophically.}

\end{frame}

\begin{frame}{PTQ vs.\ QAT: When Precision Matters}
\note{[3 min] Two strategies for quantization. PTQ is simple --- quantize the
trained model and deploy. QAT inserts fake quantization during training so
the model learns to compensate. PTQ is the strong baseline (always try first);
QAT recovers 0.5--1.0 pp when PTQ is not enough. Walk through the QAT loop
diagram: forward pass with fake quantize, backward pass with STE.
Common error: students think QAT is always better --- PTQ is often sufficient.}

\small
\begin{columns}[T]
  \begin{column}{0.48\textwidth}
    \safeimg[width=\textwidth,height=4.5cm,keepaspectratio]{qat-training-loop.pdf}
  \end{column}
  \begin{column}{0.48\textwidth}
    {\scriptsize
    \textbf{Post-Training Quantization (PTQ):}
    \begin{itemize}\setlength\itemsep{-1pt}
      \item Quantize once after training
      \item No retraining needed --- strong baseline
    \end{itemize}

    \vspace{0.05cm}
    \textbf{Quantization-Aware Training (QAT):}
    \begin{itemize}\setlength\itemsep{-1pt}
      \item Fake quantize in forward pass
      \item STE enables gradient flow
      \item Recovers 0.5--1.0 pp accuracy
    \end{itemize}
    }

    \vspace{0.05cm}
    \begin{mlsyscard}{routingstroke}
    {\scriptsize \textbf{LLMs:} INT4 weights + FP16 activations dominates (bandwidth-bound).}
    \end{mlsyscard}
  \end{column}
\end{columns}

\end{frame}

% --- ACTIVE LEARNING 2: Exercise ---
\begin{frame}{Your Turn: The Quantization Speedup}
\note{[3 min] Give students 90 seconds. FP16: 14/50 = 280 ms = 3.6 tok/s.
INT4: 3.5/50 = 70 ms = 14 tok/s. 4x speedup. Bandwidth-bound = linear.}

\footnotesize
\begin{columns}[T]
  \begin{column}{0.55\textwidth}
    {\small\bfseries Calculate the Quantization Speedup}

    \vspace{0.1cm}
    7B LLM on device with 16 GB RAM, 50 GB/s BW.\\
    FP16: 2 bytes/param. INT4: 0.5 bytes/param.

    \vspace{0.1cm}
    \textbf{Calculate:}
    \begin{enumerate}\setlength\itemsep{-1pt}
      \item FP16 model size and latency/token
      \item INT4 model size and latency/token
      \item Tokens/sec and speedup
    \end{enumerate}

    {\scriptsize\textcolor{midgray}{(90 seconds --- compare with a neighbor)}}
  \end{column}
  \begin{column}{0.42\textwidth}
    \pause
    \begin{mlsyscard}{computestroke}
    {\scriptsize \textbf{Solution:}\\[0.05cm]
    FP16: $7\!\times\!10^9\!\times\!2$ = 14 GB\\
    Latency: $14/50$ = \textbf{280 ms} (3.6 tok/s)\\[0.05cm]
    INT4: $7\!\times\!10^9\!\times\!0.5$ = 3.5 GB\\
    Latency: $3.5/50$ = \textbf{70 ms} (14 tok/s)\\[0.05cm]
    \textcolor{computestroke}{\textbf{4$\times$ linear speedup}}\\
    (bandwidth-bound $\Rightarrow$ linear)}
    \end{mlsyscard}
  \end{column}
\end{columns}

\end{frame}

% =============================================================================
\section{NAS \& Fusion}
% =============================================================================

\begin{frame}{Neural Architecture Search (NAS)}
\note{[3 min] NAS automates what human engineers do by trial and error. Four
components form a loop: search space, strategy, evaluation, and the discovered
architecture. The original NAS cost 22,400 GPU-hours. Weight-sharing reduced
this by 1000x, making NAS practical. EfficientNet: 8.4x smaller and 6.1x
faster than ResNet-50 at the same accuracy. Hardware-aware NAS includes
latency on the target device as a search objective.
If short: focus on the EfficientNet result.}

% --- Layout: FULL-WIDTH diagram ---
\centering
\safeimg[width=0.85\textwidth,height=4.2cm,keepaspectratio]{nas-overview.pdf}

\vspace{0.05cm}
\footnotesize
\begin{columns}[T]
  \begin{column}{0.48\textwidth}
    {\scriptsize \textbf{Original NAS} (2017): 22,400 GPU-hours\\
    \textbf{Weight-sharing}: 1,000$\times$ cheaper\\
    \textbf{Hardware-aware}: latency on target device}
  \end{column}
  \begin{column}{0.48\textwidth}
    \begin{mlsyscard}{datastroke}
    {\scriptsize \textbf{EfficientNet}: 8.4$\times$ smaller, 6.1$\times$ faster than ResNet-50 at same accuracy.}
    \end{mlsyscard}
  \end{column}
\end{columns}
\mlsyscite{Zoph \& Le, 2017; Tan \& Le, 2019}

\end{frame}

\begin{frame}{Operator Fusion: The Free Lunch}
\note{[3 min] This is Region 1 --- zero accuracy cost. Walk through the
left side: Conv writes to global memory, BatchNorm reads it back, writes
again, ReLU reads again. Three round-trips. The fused kernel does everything
in on-chip registers. Result: 10--30\% latency reduction for free. This is
always the first optimization to apply. Conv-BN-ReLU fusion appears in
every modern CNN inference engine. Ask: ``Why doesn't everyone do this
already?'' Answer: inference engines do --- TensorRT, ONNX Runtime, TFLite.}

% --- Layout: FULL-WIDTH diagram ---
\centering
\safeimg[width=0.88\textwidth,height=4.5cm,keepaspectratio]{operator-fusion.pdf}

\vspace{0.05cm}
\mlsysinsight{Rule:}{Operator fusion is \textbf{always} first. Zero accuracy cost, immediate latency gain.}

\end{frame}

% =============================================================================
\section{Composing Techniques}
% =============================================================================

\begin{frame}{The Decision Framework}
\note{[3 min] Map constraints to techniques. The order matters: fuse first
(free), quantize (cheap), prune (tuning), distill (training budget).}

\footnotesize
\renewcommand{\arraystretch}{1.05}
\begin{tabular}{@{}lp{2.6cm}p{2.2cm}p{2.8cm}@{}}
  \toprule
  \textbf{Constraint} & \textbf{Model Repr.} & \textbf{Numerics} & \textbf{HW Impl.} \\
  \midrule
  Memory    & Pruning, distill.\ & INT8/INT4  & --- \\
  Latency   & Channel pruning    & FP16, INT8 & Fusion, graph opt. \\
  Power     & NAS, pruning       & INT4       & Sparsity exploit. \\
  \bottomrule
\end{tabular}

\vspace{0.1cm}
\begin{mlsyscard}{crimson}
{\scriptsize \textbf{Recommended Order:} 1.\ Fusion (free) $\to$ 2.\ INT8 quant (cheap) $\to$ 3.\ Structured pruning (tuning) $\to$ 4.\ Distillation (training)}
\end{mlsyscard}

\vspace{0.05cm}
{\scriptsize \textbf{BERT}: 440 MB $\to$ 110 MB (INT8) $\to$ 28 MB (pruned + quantized + distilled) = \textbf{16$\times$} total.}

\end{frame}

\begin{frame}{The 4$\times$ MobileNet Win}
\note{[2 min] Concrete example. MobileNetV3 at 8 FPS. INT8 gets 35 FPS.
Without compression, the feature could not ship.}

\footnotesize
\begin{columns}[T]
  \begin{column}{0.55\textwidth}
    \textbf{Scenario:} Video call background blur, 30 FPS needed.

    \vspace{0.1cm}
    \renewcommand{\arraystretch}{1.1}
    \begin{tabular}{@{}lrr@{}}
      \toprule
      & \textbf{Before} & \textbf{After INT8} \\
      \midrule
      Speed     & 8 FPS           & \textbf{35 FPS} \\
      Energy    & 5 mJ/frame      & \textbf{1.7 mJ/frame} \\
      Ships?    & \textcolor{errorstroke}{No} & \textcolor{datastroke}{\textbf{Yes}} \\
      \bottomrule
    \end{tabular}

    \vspace{0.1cm}
    \mlsysconcept{Key:}{Compression \textbf{enabled} the feature.}
  \end{column}
  \begin{column}{0.42\textwidth}
    \begin{mlsyscard}{routingstroke}
    {\scriptsize \textbf{Conservation of Complexity}\\[0.05cm]
    Complexity is relocated, not destroyed:\\
    $\bullet$ Pruning $\to$ HW handles sparsity\\
    $\bullet$ Distillation $\to$ training budget\\
    $\bullet$ NAS $\to$ search compute\\
    $\bullet$ Quantization $\to$ precision lost\\[0.05cm]
    Move complexity where \textbf{cost is lowest}.}
    \end{mlsyscard}
  \end{column}
\end{columns}

\end{frame}

% --- ACTIVE LEARNING 3: Discussion ---
\begin{frame}{Discussion: Design a Compression Pipeline}
\note{[3 min] Turn-and-talk. Students discuss in pairs for 90 seconds.
Cold-call 2--3 pairs. Good answers: (1) Fuse operators, (2) INT8 quantize,
(3) structured channel pruning to fit 384 KB, (4) maybe distill from larger
model. The key constraint is 384 KB SRAM --- even MobileNetV2 INT8 at 3.5 MB
is 10x too large. They need DS-CNN or similar.}

\centering
\vspace{0.8cm}
{\large\bfseries Turn and Talk \textcolor{midgray}{(90 seconds)}}

\vspace{0.5cm}
{\normalsize You must deploy a keyword spotter (``Hey Siri'')\\
on an Arduino Nicla (\textbf{384 KB} SRAM, \textbf{100 mW}).\\[0.2cm]
\alert{Design a compression pipeline.}}

\vspace{0.4cm}
\begin{columns}[c]
  \begin{column}{0.20\textwidth}\centering\small Prune?\end{column}
  \begin{column}{0.20\textwidth}\centering\small Quantize?\end{column}
  \begin{column}{0.20\textwidth}\centering\small Distill?\end{column}
  \begin{column}{0.20\textwidth}\centering\small NAS?\end{column}
  \begin{column}{0.20\textwidth}\centering\small All?\end{column}
\end{columns}

\end{frame}

% =============================================================================
\section{Wrap-Up}
% =============================================================================

\begin{frame}{Fallacies}
\note{[2 min] Three common misconceptions. The interaction fallacy is the most
important: pruning + quantization does NOT give the product of their individual
speedups. The threshold fallacy: INT8 loses 1\% does not mean INT4 loses 2\%.
Binary ResNet-50 drops to 51\% --- catastrophic.}

\footnotesize
\textbf{Fallacy:} \textit{Optimization techniques compose independently.}\\
{\scriptsize BERT pruned to 70\%: 97.8\% accuracy. Add INT8: drops to 94.2\%. Techniques interact non-linearly.}

\vspace{0.1cm}
\textbf{Fallacy:} \textit{Aggressive quantization scales uniformly.}\\
{\scriptsize ResNet-50 INT8: 76.1\% (vs.\ 76.2\% FP32). INT4: 5--15\% drop. Binary: \textbf{51\%}. Threshold effects.}

\vspace{0.1cm}
\textbf{Fallacy:} \textit{Compression ratios predict deployment speedup.}\\
{\scriptsize INT8 + 50\% pruning $\to$ expected 8$\times$. Actual: \textbf{28\%} end-to-end. Sparsity + dequantization overhead erode gains.}

\end{frame}

\begin{frame}{Pitfalls}
\note{[2 min] Two operational pitfalls. First: optimizing for FLOPs instead of
wall-clock latency. A 60\% FLOP reduction may give only 12\% speedup on ARM
because sparse patterns prevent vectorization. Second: using PTQ when QAT
would recover 2--4\% accuracy that makes the difference between shipping and
not shipping. If short: cover just the first pitfall.}

\small
\textbf{Pitfall:} \textit{Optimizing for theoretical metrics rather than deployment performance.}\\
{\footnotesize Teams reduce FLOPs by 60\% but achieve only 12\% latency reduction on ARM. Memory bandwidth, cache behavior, and instruction-level parallelism determine actual performance --- not operation counts. \textbf{Profile on target hardware.}}

\vspace{0.2cm}
\textbf{Pitfall:} \textit{Using post-training optimization without considering training-aware alternatives.}\\
{\footnotesize PTQ achieves 96.8\% of BERT baseline. QAT retains 99.1\% --- a 2.3-point gap. Post-training pruning to 70\%: 73.8\% accuracy. Gradual pruning \emph{during} training: 75.6\%. The 2--4\% improvement from training-aware methods often determines whether models meet production thresholds.}

\end{frame}

% --- RETRIEVAL PRACTICE ---
\begin{frame}{What Were the Key Ideas?}
\note{[2 min] Retrieval practice. Students write for 90 seconds without notes.
Walk around the room. Do NOT reveal the next slide yet.}

\centering
\vspace{1.5cm}
{\Large\bfseries Close your notes.}

\vspace{0.8cm}
{\large Write down the \textbf{4 most important concepts} from today.}

\vspace{0.8cm}
{\normalsize\textcolor{midgray}{90 seconds --- no peeking.}}

\end{frame}

\begin{frame}{Key Takeaways}
\note{[2 min] Reveal takeaways. Walk through each bullet. Emphasize the
quantitative anchors: 4x from INT8, 16x combined, profile on hardware.}

\footnotesize
\begin{itemize}\setlength\itemsep{2pt}
  \item \textbf{Three-layer optimization stack}: Structural (what to compute), Precision (how precisely), Architectural (how efficiently). Combine all three for maximum compression.
  \item \textbf{PTQ is a strong baseline}: INT8 post-training quantization delivers 4$\times$ compression with no retraining. Use QAT only when baseline accuracy is insufficient.
  \item \textbf{For LLMs, weight-only quantization wins}: INT4 weights + FP16 activations dominate because generation is memory-bandwidth bound, not compute bound.
  \item \textbf{Structured pruning for commodity hardware}: Unstructured sparsity needs sparse kernels. Structured pruning (channels, heads) delivers real latency gains on GPUs.
  \item \textbf{Order matters}: Fuse first (free), quantize second (cheap), prune third (requires tuning), distill last (requires training).
  \item \textbf{Profile on target hardware}: FLOPs and parameter counts mispredict real-world performance. A 2$\times$ FLOP reduction may yield only 1.2$\times$ speedup.
\end{itemize}

\end{frame}

\begin{frame}{References}
\note{[1 min] Point students to the canonical papers.}

\small
\mlsysref{Han+15}{Han et al. ``Deep Compression.'' ICLR 2016.}
\mlsysref{Hinton+15}{Hinton, Vinyals, Dean. ``Distilling the Knowledge in a Neural Network.'' 2015.}
\mlsysref{Jacob+18}{Jacob et al. ``Quantization and Training of Neural Networks.'' CVPR 2018.}
\mlsysref{Frankle+19}{Frankle \& Carlin. ``The Lottery Ticket Hypothesis.'' ICLR 2019.}
\mlsysref{Zoph+17}{Zoph \& Le. ``Neural Architecture Search with RL.'' ICLR 2017.}
\mlsysref{Tan+19}{Tan \& Le. ``EfficientNet: Rethinking Model Scaling.'' ICML 2019.}
\mlsysref{Sanh+19}{Sanh et al. ``DistilBERT.'' 2019.}

\end{frame}

\begin{frame}{Next Lecture: Hardware Acceleration}
\note{[1 min] Forward hook. We have compressed the model's logic --- now it must
meet silicon. How do GPUs, TPUs, and NPUs exploit these optimizations?
Tensor Cores, systolic arrays, and the roofline model explain why some
compressions translate to speedup and others do not. The next chapter
completes the picture.}

\footnotesize
\begin{columns}[T]
  \begin{column}{0.55\textwidth}
    \vspace{0.3cm}
    We compressed the model's \textbf{logic} --- shaved off every unnecessary bit.

    \vspace{0.2cm}
    Logic must eventually run on \textbf{physics}.

    \vspace{0.2cm}
    \begin{mlsyscard}{crimson}
    \textbf{Next:} How do accelerators exploit compression?\\[0.1cm]
    {\footnotesize
    $\bullet$ Tensor Cores for INT8/INT4 arithmetic\\
    $\bullet$ Systolic arrays for dense matrix ops\\
    $\bullet$ Roofline model: when does compression help?\\
    $\bullet$ Sparse Tensor Cores for 2:4 sparsity}
    \end{mlsyscard}
  \end{column}
  \begin{column}{0.42\textwidth}
    \vspace{0.5cm}
    \centering
    {\normalsize\textbf{From Math to Physics}}

    \vspace{0.3cm}
    {\large Compressed model $+$\\
    Smart hardware $=$\\
    \textcolor{datastroke}{\textbf{Real speedup}}}

    \vspace{0.3cm}
    {\footnotesize Without hardware alignment,\\
    compression ratios are just\\
    numbers on paper.}
  \end{column}
\end{columns}

\end{frame}

\end{document}
